\section{The Crystal: The Substrate}
\label{sec:crystal}

Now the picture made concrete. The substrate is a hyperbolic crystal, and the simplest face
of it can be drawn. Figure~\ref{fig:mesh} shows the heptagrid, the regular tiling of the
hyperbolic plane by seven-sided cells, three meeting at each corner, written in the Schl\"afli
shorthand as $\{7,3\}$.

\begin{figure}[h]
\centering
\includegraphics[width=0.5\textwidth]{vibe-mesh-73.pdf}
\caption{The $\{7,3\}$ crystal, the simplest face of the substrate. Each tile is a vibe.
Its color is its tone: red is pain, green is peace, blue is pleasure. Tiles that touch are
vibes that note (experience) one another, so the edges are the relations of the mesh.}
\label{fig:mesh}
\end{figure}

The picture is meant literally. Each tile is a \emph{vibe}, one unit of experience. Its color
is its \emph{tone}, its felt charge: red is pain, green is peace, blue is pleasure. Two tiles
that touch are two vibes that note, that is, experience, one another, so the edges of the
crystal are the notes of the mesh. There is nothing else in the model but this: colored tiles
and the touching between them. The tessellation drawn in tone-colors is not an artist's
impression. It is intended, within the model, as a literal map of the substance.

\paragraph{Why hyperbolic.}
The geometry is not a free choice. It is forced by two demands. First, the universe grows
without end, and a flat space has too little room: the number of cells within a given distance
grows only polynomially, too slowly to hold an endlessly branching, eternally growing
structure. Hyperbolic space opens out exponentially, so there is always more room at the
edge. Second, and more importantly, a flat regular grid has preferred directions, along its rows
and columns, which would tend to make the speed of light depend on direction and so strain the
relativity that experiment confirms to high precision. A hyperbolic crystal does not have this
problem.
In curved space there is no global notion of ``the same direction'' at two different places,
so the curvature scrambles the directions and a regular hyperbolic tiling reads, to any local
observer, much like a random foam with no preferred frame. What the testbed suggests, then, is
that \emph{curvature, rather than disorder, may be what buys Lorentz safety}. On this picture a
universe could be an ordered crystal and still look smooth and relativistic from within.

\paragraph{Why three dimensions.}
The flat $\{7,3\}$ picture is the easy-to-draw special case. The real substrate is one member
of a whole family of hyperbolic crystals, the regular honeycombs of a reflection group, and
the actual universe is the three-dimensional one, the dodecahedral honeycomb $\{5,3,4\}$, a
crystal of twelve-sided cells filling curved three-dimensional space, with time as its growth.
The dimension is not assumed. It is constrained twice over. A crystal with finite cells, which
is what a substrate of finite, locally connected vibes must be, exists only in two, three, and
four dimensions: above four there is no way to fill hyperbolic space with finite cells at all.
And within that narrow window, only three spatial dimensions supports stable, closed
gravitational orbits, the inverse-square law that lets atoms and solar systems hold together,
while two precesses and four is unstable. So the universe we can live in is three-dimensional
in space, four-dimensional with time. The two-dimensional tiling is shown here only because it
can be drawn on a page. The model is robust across the family, and a little richer and harder
to picture than the figure suggests.

\paragraph{One object, many views.}
It is worth holding all of these together, because they are not separate claims but one object
seen from different angles. As numbers, the substrate is the integers. As algebra, it is their
symmetry group, generalized to reflections. As geometry, it is a hyperbolic crystal. As a
felt thing, it is a web of vibes. The integers are why the crystal is forced rather than
chosen, the geometry is its shape, and the vibes are what actually sit at every tile. The
figure is the relations drawn so the shape is visible, not a container the vibes live in.
